\chapter{Chapter08 — Recursive NAND Composition}
%%% generated by the blueprint pipeline from FABL.Chapter08.RecursiveNANDComposition
\section{Overview}

Declarations of the Lean module \texttt{FABL.Chapter08.RecursiveNANDComposition}.

\section{Declarations}

\begin{definition}[Decision tree reindexing]
\lean{FABL.FiniteDecisionTree.castAvailable}
\label{FABL.FiniteDecisionTree.castAvailable}
\leanok
\uses{FABL.FiniteDecisionTree}
For sets of available coordinates $A$ and $B$ and an equality $A=B$, the reindexing
operation views any finite decision tree whose available coordinates are $A$ as a tree
whose available coordinates are $B$.
\end{definition}

\begin{theorem}[Evaluation under availability-set casting]
\lean{FABL.FiniteDecisionTree.eval_castAvailable}
\label{FABL.FiniteDecisionTree.eval_castAvailable}
\leanok
\uses{FABL.FiniteDecisionTree, FABL.FiniteDecisionTree.castAvailable, FABL.FiniteDecisionTree.eval}
\difficulty{1}
Let $T$ be a finite decision tree whose available coordinates are in a finite set $A$,
and let $B$ be another such set with $A=B$. For every input $x$, casting $T$ to have
available set $B$ leaves its evaluation unchanged.
\end{theorem}

\begin{theorem}[Query trace under availability identification]
\lean{FABL.FiniteDecisionTree.queryTrace_castAvailable}
\label{FABL.FiniteDecisionTree.queryTrace_castAvailable}
\leanok
\uses{FABL.FiniteDecisionTree, FABL.FiniteDecisionTree.castAvailable, FABL.FiniteDecisionTree.queryTrace}
\difficulty{1}
Let $n$ be a nonnegative integer, let $\Omega$ be a set of coordinate values, and let
$\alpha$ be a set of leaf labels. A finite decision tree queries available coordinates
and branches according to their observed values. If $A$ and $B$ are equal sets of
coordinates, then identifying a tree whose available set is $A$ as a tree whose
available set is $B$ leaves its query trace unchanged on every input $x$.
\end{theorem}

\begin{theorem}[Query count under a change of available set]
\lean{FABL.FiniteDecisionTree.queryCount_castAvailable}
\label{FABL.FiniteDecisionTree.queryCount_castAvailable}
\leanok
\uses{FABL.FiniteDecisionTree, FABL.FiniteDecisionTree.castAvailable, FABL.FiniteDecisionTree.queryCount, FABL.FiniteDecisionTree.queryTrace_castAvailable}
\difficulty{3}
Let $n$ be a natural number, let $\Omega$ and $\alpha$ be types, and let
$A,B\subseteq\{0,\ldots,n-1\}$ be equal sets. For any finite decision tree $T$ with
available-coordinate set $A$ and any input $x:\{0,\ldots,n-1\}\to\Omega$, regarding $T$
as a tree with available set $B$ leaves its query count on $x$ unchanged.
\end{theorem}

\begin{definition}[Decision tree weakening]
\lean{FABL.FiniteDecisionTree.weaken}
\label{FABL.FiniteDecisionTree.weaken}
\leanok
\uses{FABL.FiniteDecisionTree}
Let $A$ and $B$ be finite sets of coordinates with $A\subseteq B$. Any finite decision
tree that queries only coordinates in $A$ can also be regarded as a tree with the larger
set $B$ of coordinates available.
\end{definition}

\begin{theorem}[Evaluation after enlarging available coordinates]
\lean{FABL.FiniteDecisionTree.eval_weaken}
\label{FABL.FiniteDecisionTree.eval_weaken}
\leanok
\uses{FABL.FiniteDecisionTree, FABL.FiniteDecisionTree.eval, FABL.FiniteDecisionTree.weaken}
\difficulty{4}
Let $A\subseteq B$ be sets of available coordinates among $n$ coordinates, and let $T$
be a finite decision tree with labels in $\alpha$ and available coordinates $A$. For any
input $x$ assigning an alphabet symbol to each coordinate, evaluating $T$ viewed as a
tree with the larger available set $B$ gives the same label as evaluating $T$ at $x$.
\end{theorem}

\begin{theorem}[Trace under enlargement of the available set]
\lean{FABL.FiniteDecisionTree.queryTrace_weaken}
\label{FABL.FiniteDecisionTree.queryTrace_weaken}
\leanok
\uses{FABL.FiniteDecisionTree, FABL.FiniteDecisionTree.queryTrace, FABL.FiniteDecisionTree.weaken}
\difficulty{4}
Let $A\subseteq B\subseteq\{0,\ldots,n-1\}$, and let $T$ be a finite decision tree whose
queries are available in $A$. Regard the same tree as having queries available in $B$.
For every input, its query trace is unchanged.
\end{theorem}

\begin{definition}[Decision tree output mapping]
\lean{FABL.FiniteDecisionTree.mapOutput}
\label{FABL.FiniteDecisionTree.mapOutput}
\leanok
\uses{FABL.FiniteDecisionTree}
Let $g$ be a map from an output set $\alpha$ to an output set $\beta$. Applying $g$ to
the output labels of a finite decision tree replaces each leaf label $a$ by $g(a)$ and
leaves the queried coordinates and branching structure unchanged.
\end{definition}

\begin{theorem}[Evaluation after relabeling outputs]
\lean{FABL.FiniteDecisionTree.eval_mapOutput}
\label{FABL.FiniteDecisionTree.eval_mapOutput}
\leanok
\uses{FABL.FiniteDecisionTree, FABL.FiniteDecisionTree.eval, FABL.FiniteDecisionTree.mapOutput}
\difficulty{4}
Let $T$ be a finite decision tree with input alphabet $\Omega$, output labels in $\alpha$, and
available coordinates $A\subseteq\{0,\ldots,n-1\}$. For any function $g$ from its
output-label set to a set $\beta$, and any input $x$ assigning an alphabet symbol to each
of the $n$ coordinates, evaluating the tree obtained by applying $g$ to every leaf label
gives $g$ applied to the evaluation of $T$ at $x$.
\end{theorem}

\begin{theorem}[Query trace under relabeling leaves]
\lean{FABL.FiniteDecisionTree.queryTrace_mapOutput}
\label{FABL.FiniteDecisionTree.queryTrace_mapOutput}
\leanok
\uses{FABL.FiniteDecisionTree, FABL.FiniteDecisionTree.mapOutput, FABL.FiniteDecisionTree.queryTrace}
\difficulty{4}
Let $n$ be a nonnegative integer, let $\Omega,\alpha,\beta$ be sets, and let
$g:\alpha\to\beta$ be a map. For any finite decision tree over $\Omega$ with leaf labels
in $\alpha$, replacing each leaf label $a$ by $g(a)$ leaves the query trace unchanged on
every input $x$.
\end{theorem}

\begin{definition}[Reindexing a decision tree]
\lean{FABL.FiniteDecisionTree.reindex}
\label{FABL.FiniteDecisionTree.reindex}
\leanok
\uses{FABL.FiniteDecisionTree, FABL.FiniteDecisionTree.castAvailable}
Given an injective map from a finite coordinate set of size $m$ into one of size $n$,
reindex a decision tree whose available coordinates form a subset $A$ of the first set
as a decision tree whose available coordinates are the image of $A$. The tree's labels
and branching choices are preserved, with each queried coordinate replaced by its image.
\end{definition}

\begin{theorem}[Evaluation after coordinate reindexing]
\lean{FABL.FiniteDecisionTree.eval_reindex}
\label{FABL.FiniteDecisionTree.eval_reindex}
\leanok
\uses{FABL.FiniteDecisionTree, FABL.FiniteDecisionTree.eval, FABL.FiniteDecisionTree.eval_castAvailable, FABL.FiniteDecisionTree.reindex}
\difficulty{4}
Let $e$ be an embedding of the $m$ coordinates into the $n$ coordinates, let $A$ be a
set of available coordinates among the $m$ coordinates, and let $T$ be a finite decision
tree with labels in $\alpha$ and available coordinates $A$. For any input $x$ assigning
an alphabet symbol to each of the $n$ coordinates, evaluating the reindexed tree equals
evaluating $T$ at the input $i\mapsto x(e(i))$.
\end{theorem}

\begin{theorem}[Query trace length under coordinate embedding]
\lean{FABL.FiniteDecisionTree.queryTrace_reindex_length}
\label{FABL.FiniteDecisionTree.queryTrace_reindex_length}
\leanok
\uses{FABL.FiniteDecisionTree, FABL.FiniteDecisionTree.queryTrace, FABL.FiniteDecisionTree.queryTrace_castAvailable, FABL.FiniteDecisionTree.reindex}
\difficulty{4}
Let $m,n$ be nonnegative integers, let $\Omega$ be a set of coordinate values, and let
$\alpha$ be a set of leaf labels. Let $e:\{0,\ldots,m-1\}\to\{0,\ldots,n-1\}$ be an
injective map, let $A$ be a finite set of coordinates among the first $m$ coordinates,
and let $T$ be a finite decision tree over $\Omega$ with leaf labels in $\alpha$ and
available coordinates $A$. For any input $x$ on the $n$ coordinates, the length of the
query trace of the tree reindexed by $e$ equals the length of the query trace of $T$ on
the input $i\mapsto x(e(i))$.
\end{theorem}

\begin{definition}[Composition on disjoint coordinates]
\lean{FABL.FiniteDecisionTree.thenOnDisjoint}
\label{FABL.FiniteDecisionTree.thenOnDisjoint}
\leanok
\uses{FABL.FiniteDecisionTree, FABL.FiniteDecisionTree.castAvailable, FABL.FiniteDecisionTree.weaken}
Let $A$ and $B$ be disjoint finite sets of coordinates, let $T$ be a decision tree over
$A$ with labels in $\alpha$, and let $K$ assign to each label in $\alpha$ a decision
tree over $B$ with labels in $\beta$. The resulting tree over $A\cup B$ runs $T$ first,
then follows the tree assigned by $K$ to the label returned by $T$.
\end{definition}

\begin{theorem}[Evaluation of sequential decision trees]
\lean{FABL.FiniteDecisionTree.eval_thenOnDisjoint}
\label{FABL.FiniteDecisionTree.eval_thenOnDisjoint}
\leanok
\uses{FABL.FiniteDecisionTree, FABL.FiniteDecisionTree.eval, FABL.FiniteDecisionTree.eval_castAvailable, FABL.FiniteDecisionTree.eval_weaken, FABL.FiniteDecisionTree.thenOnDisjoint}
\difficulty{4}
Let $A$ and $B$ be disjoint sets of coordinates among $n$ coordinates. Let $T$ be a
finite decision tree with labels in $\alpha$ and available coordinates $A$, and let $K$
assign to each label in $\alpha$ a finite decision tree with labels in $\beta$ and
available coordinates $B$. For every input $x$ assigning an alphabet symbol to each
coordinate, evaluating the tree formed by running $T$ and then the tree $K$ selected by
its output equals evaluating $K(T(x))$ at the same input $x$.
\end{theorem}

\begin{theorem}[Trace of a sequentially composed tree]
\lean{FABL.FiniteDecisionTree.queryTrace_thenOnDisjoint}
\label{FABL.FiniteDecisionTree.queryTrace_thenOnDisjoint}
\leanok
\uses{FABL.FiniteDecisionTree, FABL.FiniteDecisionTree.eval, FABL.FiniteDecisionTree.queryTrace, FABL.FiniteDecisionTree.queryTrace_castAvailable, FABL.FiniteDecisionTree.queryTrace_weaken, FABL.FiniteDecisionTree.thenOnDisjoint}
\difficulty{4}
Let $A,B\subseteq\{0,\ldots,n-1\}$ be disjoint, let $T$ be a finite decision tree
querying coordinates in $A$, and let $K$ assign to each output of $T$ a finite decision
tree querying coordinates in $B$. For every input $x$, the query trace of the tree that
runs $T$ and then runs $K$ on the output of $T$ is the concatenation of the trace of $T$
and the trace of $K(T(x))$.
\end{theorem}

\begin{theorem}[Query count of sequential disjoint trees]
\lean{FABL.FiniteDecisionTree.queryCount_thenOnDisjoint}
\label{FABL.FiniteDecisionTree.queryCount_thenOnDisjoint}
\leanok
\uses{FABL.FiniteDecisionTree, FABL.FiniteDecisionTree.eval, FABL.FiniteDecisionTree.queryCount, FABL.FiniteDecisionTree.queryTrace_thenOnDisjoint, FABL.FiniteDecisionTree.thenOnDisjoint}
\difficulty{3}
Let $n$ be a nonnegative integer, let $\Omega$ be a set of coordinate values, and let
$\alpha,\beta$ be sets of leaf labels. Consider finite decision trees that query only
coordinates in their specified available sets and branch according to observed values.
Let $A,B\subseteq\{0,\ldots,n-1\}$ be disjoint, let $T$ be a tree using coordinates in
$A$, and for each $a\in\alpha$ let $K(a)$ be a tree using coordinates in $B$. For any
input $x$, the number of queries in the tree that first runs $T$ and then runs $K$ on
the resulting label is the number of queries in $T$ plus the number in $K(T(x))$, each
evaluated on $x$.
\end{theorem}

\begin{definition}[Recursive NAND block embedding]
\lean{FABL.recursiveNandBlockEmbedding}
\label{FABL.recursiveNandBlockEmbedding}
\leanok
For each natural number $d$ and each $i\in\{0,1\}$, define an injective map from the
coordinates in $\mathrm{Fin}(2^d)$ into those in $\mathrm{Fin}(2^{d+1})$ by sending $j$
to the coordinate corresponding to the pair $(i,j)$ under the standard identification of
$\mathrm{Fin}(2^{d+1})$ with $\mathrm{Fin}(2)\times\mathrm{Fin}(2^d)$.
\end{definition}

\begin{definition}[Recursive NAND block coordinates]
\lean{FABL.recursiveNandBlockCoordinates}
\label{FABL.recursiveNandBlockCoordinates}
\leanok
\uses{FABL.recursiveNandBlockEmbedding}
For each natural number $d$ and each $i\in\{0,1\}$, the recursive NAND block coordinates
are the image, under the recursive NAND block embedding, of all coordinates in
$\mathrm{Fin}(2^d)$.
\end{definition}

\begin{theorem}[Disjoint recursive NAND blocks]
\lean{FABL.recursiveNandBlockCoordinates_disjoint}
\label{FABL.recursiveNandBlockCoordinates_disjoint}
\leanok
\uses{FABL.recursiveNandBlockCoordinates}
\difficulty{4}
For every natural number $d$, the two sets of coordinates occupied by the two blocks of
size $2^d$ in a vector with $2^{d+1}$ coordinates are disjoint.
\end{theorem}

\begin{theorem}[Union of recursive block coordinates]
\lean{FABL.recursiveNandBlockCoordinates_union}
\label{FABL.recursiveNandBlockCoordinates_union}
\leanok
\uses{FABL.recursiveNandBlockCoordinates, FABL.recursiveNandBlockEmbedding}
\difficulty{4}
For every natural number $d$, the coordinates in the two blocks of size $2^d$ that
partition a block of size $2^{d+1}$ have union equal to the full set of $2^{d+1}$
coordinates.
\end{theorem}

\begin{definition}[NAND of two signs]
\lean{FABL.nandPair}
\label{FABL.nandPair}
\leanok
\uses{FABL.Sign}
For signs $x,y\in\{\pm1\}$, their NAND is $+1$ if both signs are $-1$, and is $-1$
otherwise.
\end{definition}

\begin{theorem}[NAND function on two inputs]
\lean{FABL.nandFunction_eq_nandPair}
\label{FABL.nandFunction_eq_nandPair}
\leanok
\uses{FABL.SignCube, FABL.nandFunction, FABL.nandPair}
\difficulty{1}
For any $x\in\{-1,1\}^2$, the two-input function that returns $1$ exactly when both
coordinates are $-1$, and $-1$ otherwise, equals the operation on the pair $(x_0,x_1)$
that returns $1$ exactly when both entries are $-1$, and $-1$ otherwise.
\end{theorem}

\begin{definition}[Recursive NAND identity tree]
\lean{FABL.recursiveNandIdentityTree}
\label{FABL.recursiveNandIdentityTree}
\leanok
\uses{FABL.FiniteAlphabetDecisionTree, FABL.Sign}
The recursive NAND identity tree on one sign coordinate queries that coordinate and
returns the observed sign as its label.
\end{definition}

\begin{theorem}[Identity tree computes the initial function]
\lean{FABL.recursiveNandIdentityTree_computes}
\label{FABL.recursiveNandIdentityTree_computes}
\leanok
\uses{FABL.FiniteDecisionTree.Computes, FABL.recursiveNAND, FABL.recursiveNandIdentityTree}
\difficulty{1}
On a single sign input, the decision tree that queries its sole coordinate and returns
the observed sign computes the depth-zero recursive NAND function.
\end{theorem}

\begin{definition}[NAND continuation tree]
\lean{FABL.nandContinuationTree}
\label{FABL.nandContinuationTree}
\leanok
\uses{FABL.FiniteDecisionTree, FABL.FiniteDecisionTree.mapOutput, FABL.Sign, FABL.nandPair}
Given a sign $x$ and a decision tree on a finite coordinate set $B$ whose labels are
signs, define a tree on $B$ that returns $-1$ immediately if $x=+1$; otherwise, it
applies the NAND operation with $x$ to each output of the given tree.
\end{definition}

\begin{theorem}[Evaluation of the NAND Continuation Tree]
\lean{FABL.eval_nandContinuationTree}
\label{FABL.eval_nandContinuationTree}
\leanok
\uses{FABL.FiniteDecisionTree, FABL.FiniteDecisionTree.eval, FABL.Sign, FABL.SignCube, FABL.nandContinuationTree, FABL.nandPair}
\difficulty{4}
Let $n$ be a nonnegative integer, let $B$ be a set of available coordinates among $n$
coordinates, let $T$ be a finite decision tree on those coordinates with sign outputs,
let $a$ be a sign, and let $x$ be a sign vector on all $n$ coordinates. Evaluating the
continuation tree formed from $a$ and $T$ gives $1$ exactly when $a=-1$ and the
evaluation of $T$ on $x$ is $-1$; otherwise it gives $-1$.
\end{theorem}

\begin{theorem}[Query count of the NAND continuation tree]
\lean{FABL.queryCount_nandContinuationTree}
\label{FABL.queryCount_nandContinuationTree}
\leanok
\uses{FABL.FiniteDecisionTree, FABL.FiniteDecisionTree.queryCount, FABL.FiniteDecisionTree.queryTrace, FABL.Sign, FABL.SignCube, FABL.nandContinuationTree}
\difficulty{4}
Let $n$ be a natural number, let $B$ be a set of available coordinates among $n$
coordinates, let $a\in\{−1,1\}$, and let $T$ be a decision tree using coordinates in
$B$. For any sign-valued input $x\in\{−1,1\}^{n}$, the NAND continuation tree makes zero
queries if $a=1$, and otherwise makes exactly as many queries as $T$ does on $x$.
\end{theorem}

\begin{definition}[Ordered NAND tree]
\lean{FABL.orderedNandTree}
\label{FABL.orderedNandTree}
\leanok
\uses{FABL.FiniteDecisionTree, FABL.FiniteDecisionTree.thenOnDisjoint, FABL.Sign, FABL.nandContinuationTree}
Let $A$ and $B$ be disjoint finite coordinate sets, and let the first and second
decision trees use coordinates in $A$ and $B$, respectively, with sign labels. The
ordered NAND tree runs the first tree, then uses its output to define a NAND
continuation of the second tree, yielding a tree on $A\cup B$.
\end{definition}

\begin{theorem}[Evaluation of the ordered NAND tree]
\lean{FABL.eval_orderedNandTree}
\label{FABL.eval_orderedNandTree}
\leanok
\uses{FABL.FiniteDecisionTree, FABL.FiniteDecisionTree.eval, FABL.Sign, FABL.SignCube, FABL.nandPair, FABL.orderedNandTree}
\difficulty{1}
Let $n$ be a natural number, and let $A$ and $B$ be disjoint sets of coordinates in
$	ext{Fin}(n)$. Given two decision trees on these coordinate sets, with inputs and
outputs in $\{ -1,1\}$, their ordered composition evaluates on every input to $-1$
exactly when both tree outputs are $-1$, and to $1$ otherwise.
\end{theorem}

\begin{theorem}[Query count of an ordered NAND tree]
\lean{FABL.queryCount_orderedNandTree}
\label{FABL.queryCount_orderedNandTree}
\leanok
\uses{FABL.FiniteDecisionTree, FABL.FiniteDecisionTree.eval, FABL.FiniteDecisionTree.queryCount, FABL.FiniteDecisionTree.queryCount_thenOnDisjoint, FABL.Sign, FABL.SignCube, FABL.orderedNandTree, FABL.queryCount_nandContinuationTree}
\difficulty{3}
Let $A$ and $B$ be disjoint sets of coordinates in a finite sign cube, and let the first
and second finite decision trees use coordinates from $A$ and $B$, respectively. On an
input $x$, the ordered NAND tree first runs the first tree and runs the second tree only
if the first tree outputs $-1$. Its number of queries is the number made by the first
tree plus zero if that tree outputs $1$, and otherwise the number made by the second
tree.
\end{theorem}

\begin{definition}[Recursive NAND seed space]
\lean{FABL.RecursiveNandSeed}
\label{FABL.RecursiveNandSeed}
\leanok
Define the seed space at depth $0$ to be a singleton. At depth $d+1$, a seed consists of
an element of $\{0,1\}$ together with an ordered pair of depth-$d$ seeds.
\end{definition}

\begin{definition}[Decision tree on a recursive block]
\lean{FABL.recursiveNandTreeInBlock}
\label{FABL.recursiveNandTreeInBlock}
\leanok
\uses{FABL.FiniteAlphabetDecisionTree, FABL.FiniteDecisionTree, FABL.FiniteDecisionTree.reindex, FABL.Sign, FABL.recursiveNandBlockCoordinates, FABL.recursiveNandBlockEmbedding}
For a depth $d$, a block index $i\in\{0,1\}$, and a finite-alphabet decision tree on
$2^d$ coordinates, reindex its coordinates into the block of $2^d$ coordinates selected
by $i$ among the $2^{d+1}$ coordinates. The resulting tree has precisely that block as
its available coordinate set.
\end{definition}

\begin{theorem}[Evaluation of a decision tree on one block]
\lean{FABL.eval_recursiveNandTreeInBlock}
\label{FABL.eval_recursiveNandTreeInBlock}
\leanok
\uses{FABL.FiniteAlphabetDecisionTree, FABL.FiniteDecisionTree.eval, FABL.FiniteDecisionTree.eval_reindex, FABL.Sign, FABL.SignCube, FABL.recursiveInputBlock, FABL.recursiveNandBlockCoordinates, FABL.recursiveNandTreeInBlock}
\difficulty{3}
Let $d$ be a nonnegative integer, let $i\in\{0,1\}$, and let $T$ be a finite decision
tree whose coordinates and input symbols are indexed by $\{0,\ldots,2^d-1\}$ and whose
outputs are signs. For every sign vector $x$ indexed by $\{0,\ldots,2^{d+1}-1\}$,
evaluating the tree obtained by embedding its coordinates into block $i$ of $x$ gives
the same output as evaluating $T$ on the restriction of $x$ to that block.
\end{theorem}

\begin{theorem}[Query count in a recursive NAND block]
\lean{FABL.queryCount_recursiveNandTreeInBlock}
\label{FABL.queryCount_recursiveNandTreeInBlock}
\leanok
\uses{FABL.FiniteAlphabetDecisionTree, FABL.FiniteDecisionTree.queryCount, FABL.FiniteDecisionTree.queryTrace_reindex_length, FABL.Sign, FABL.SignCube, FABL.recursiveInputBlock, FABL.recursiveNandBlockCoordinates, FABL.recursiveNandTreeInBlock}
\difficulty{3}
Let $d$ be a nonnegative integer, let $i\in\{0,1\}$, and let $T$ be a finite decision
tree on sign inputs with $2^d$ coordinates. Embed those coordinates as block $i$ among
the $2^{d+1}$ coordinates, and let $x$ be a sign assignment on the larger cube. The
number of queries made by the reindexed tree on $x$ equals the number made by $T$ on the
restriction of $x$ to block $i$.
\end{theorem}

\begin{definition}[Recursive NAND decision tree]
\lean{FABL.recursiveNandTree}
\statementsource{Analysis-of-Boolean-Functions-by-Ryan-ODonnell}{pdf-ffbe9c7b3306-p242-b005}
\label{FABL.recursiveNandTree}
\leanok
\uses{FABL.FiniteAlphabetDecisionTree, FABL.FiniteDecisionTree.castAvailable, FABL.RecursiveNandSeed, FABL.Sign, FABL.orderedNandTree, FABL.recursiveNandBlockCoordinates_disjoint, FABL.recursiveNandBlockCoordinates_union, FABL.recursiveNandIdentityTree, FABL.recursiveNandTreeInBlock}
Define a decision tree for each depth $d$ and recursive NAND seed at depth $d$, on $2^d$
sign coordinates. At depth zero it queries the sole coordinate and returns its sign; at
each subsequent depth it recursively builds trees on the two halves, combines them in
the order selected by the seed, and computes their NAND.
\end{definition}

\begin{theorem}[Correctness of the recursive decision tree]
\lean{FABL.recursiveNandTree_computes}
\label{FABL.recursiveNandTree_computes}
\leanok
\uses{FABL.FiniteDecisionTree.Computes, FABL.FiniteDecisionTree.eval_castAvailable, FABL.RecursiveNandSeed, FABL.eval_orderedNandTree, FABL.eval_recursiveNandTreeInBlock, FABL.nandFunction, FABL.nandPair, FABL.recursiveInputBlock, FABL.recursiveNAND, FABL.recursiveNandIdentityTree_computes, FABL.recursiveNandTree}
\difficulty{5}
For every $d\in\mathbb{N}$ and every random seed at level $d$, the resulting
deterministic decision tree computes the recursively defined NAND function on inputs in
$\{\pm1\}^{2^d}$.
\end{theorem}

\begin{theorem}[Recursive NAND tree query count]
\lean{FABL.queryCount_recursiveNandTree_succ}
\label{FABL.queryCount_recursiveNandTree_succ}
\leanok
\uses{FABL.FiniteDecisionTree.queryCount, FABL.FiniteDecisionTree.queryCount_castAvailable, FABL.RecursiveNandSeed, FABL.SignCube, FABL.eval_recursiveNandTreeInBlock, FABL.queryCount_orderedNandTree, FABL.queryCount_recursiveNandTreeInBlock, FABL.recursiveInputBlock, FABL.recursiveNAND, FABL.recursiveNandTree, FABL.recursiveNandTree_computes}
\difficulty{4}
Let $d$ be a nonnegative integer, let $o\in\{0,1\}$, choose left and right seeds for
depth $d$, and let $x$ be a sign assignment on $2^{d+1}$ coordinates. If $o=0$, the
recursive NAND tree's query count is the left tree's query count on block $0$, plus zero
when the recursive NAND value on block $0$ is $1$ and otherwise the right tree's query
count on block $1$. If $o=1$, the same rule applies with the right tree on block $1$
first and the left tree on block $0$ as the possible continuation.
\end{theorem}

\begin{definition}[Recursive NAND seed distribution]
\lean{FABL.recursiveNandSeedLaw}
\label{FABL.recursiveNandSeedLaw}
\leanok
\uses{FABL.RecursiveNandSeed, FABL.uniformPMF}
Define a probability distribution on seeds recursively: at depth $0$ it is concentrated
on the unique seed. At depth $d+1$, choose the order uniformly from $\{0,1\}$, choose
two depth-$d$ seeds independently according to the depth-$d$ seed distribution, and
combine them with that order.
\end{definition}

\begin{definition}[Randomized recursive NAND tree]
\lean{FABL.recursiveNandRandomizedTree}
\statementsource{Analysis-of-Boolean-Functions-by-Ryan-ODonnell}{pdf-ffbe9c7b3306-p242-b005}
\label{FABL.recursiveNandRandomizedTree}
\leanok
\uses{FABL.RandomizedDecisionTree, FABL.RecursiveNandSeed, FABL.Sign, FABL.recursiveNAND, FABL.recursiveNandSeedLaw, FABL.recursiveNandTree, FABL.recursiveNandTree_computes}
For each nonnegative integer $d$, define a randomized decision tree on $2^d$ sign
coordinates that computes the depth-$d$ recursive NAND function exactly. Its seed space
is the recursively defined depth-$d$ NAND seed space, equipped with its specified
probability distribution, and each seed selects the corresponding deterministic
recursive NAND tree.
\end{definition}

\begin{theorem}[Adding a constant inside expectation]
\lean{FABL.pmfExpectation_const_add_general}
\label{FABL.pmfExpectation_const_add_general}
\leanok
\uses{FABL.pmfExpectation, FABL.pmfExpectation_add_fun, FABL.pmfExpectation_const}
\difficulty{3}
Let $\Omega$ be a finite set, let $\pi$ be a probability mass function on $\Omega$, let
$c\in\bbr$, and let $f:\Omega\to\bbr$. The expectation of the function $x\mapsto c+f(x)$
equals $c$ plus the expectation of $f$.
\end{theorem}

\begin{theorem}[Expectation of a constant shift]
\lean{FABL.pmfExpectation_add_const_general}
\label{FABL.pmfExpectation_add_const_general}
\leanok
\uses{FABL.pmfExpectation, FABL.pmfExpectation_add_fun, FABL.pmfExpectation_const}
\difficulty{3}
Let $\Omega$ be a finite set, let $\pi$ be a probability mass function on $\Omega$, let
$f:\Omega\to\mathbb R$, and let $c\in\mathbb R$. Then $\mathbb E_\pi(f+c)=\mathbb E_\pi
f+c$.
\end{theorem}

\begin{theorem}[Recursive NAND input-cost recurrence]
\lean{FABL.recursiveNandRandomizedTree_inputCost_succ}
\label{FABL.recursiveNandRandomizedTree_inputCost_succ}
\leanok
\uses{FABL.FiniteDecisionTree.queryCount, FABL.RandomizedDecisionTree.inputCost, FABL.SignCube, FABL.pmfExpectation, FABL.pmfExpectation_add_const_general, FABL.pmfExpectation_bind, FABL.pmfExpectation_const, FABL.pmfExpectation_const_add_general, FABL.pmfExpectation_map, FABL.pmfExpectation_uniformPMF_eq_expect, FABL.queryCount_recursiveNandTree_succ, FABL.recursiveInputBlock, FABL.recursiveNAND, FABL.recursiveNandRandomizedTree, FABL.recursiveNandSeedLaw, FABL.recursiveNandTree}
\difficulty{5}
Let $d\in\mathbb{N}$ and let $x\in\{\pm1\}^{2^{d+1}}$. Write $x^{(0)}$ and $x^{(1)}$ for
the two consecutive blocks of $2^d$ coordinates of $x$. For each $k$, let $F_k$ be the
recursively defined NAND function on $2^k$ signs, where $F_0(y)=y$ and $F_{k+1}(y)$ is
$1$ if both block values $F_k(y^{(0)})$ and $F_k(y^{(1)})$ are $-1$, and is $-1$
otherwise; let $C_k(y)$ be the expected number of queries made on input $y$ by the
randomized tree at level $k$. Then
\[
C_{d+1}(x)=\frac12\Bigl(C_d(x^{(0)})+\begin{cases}0,&F_d(x^{(0)})=1,\\
C_d(x^{(1)}),&F_d(x^{(0)})\ne1,\end{cases}+C_d(x^{(1)})+\begin{cases}0,&F_d(x^{(1)})=1,\\
C_d(x^{(0)}),&F_d(x^{(1)})\ne1.\end{cases}\Bigr).
\]
\end{theorem}

\begin{definition}[Output-conditional decision tree cost]
\lean{FABL.RandomizedDecisionTree.outputConditionalCost}
\label{FABL.RandomizedDecisionTree.outputConditionalCost}
\leanok
\uses{FABL.RandomizedDecisionTree, FABL.RandomizedDecisionTree.inputCost}
For a randomized decision tree computing $f$ on a finite, nonempty input alphabet, and
for an output value $v$, define its output-conditional cost to be the maximum, over all
inputs, of the input cost when $f$ equals $v$, and zero otherwise.
\end{definition}

\begin{theorem}[Nonnegativity of conditional query cost]
\lean{FABL.RandomizedDecisionTree.outputConditionalCost_nonneg}
\label{FABL.RandomizedDecisionTree.outputConditionalCost_nonneg}
\leanok
\uses{FABL.RandomizedDecisionTree, FABL.RandomizedDecisionTree.inputCost, FABL.RandomizedDecisionTree.inputCost_nonneg, FABL.RandomizedDecisionTree.outputConditionalCost}
\difficulty{4}
Let $\Omega$ be a finite nonempty input alphabet, let $\alpha$ be an output type with
decidable equality, and let $f:\Omega^n\to\alpha$. For a randomized decision tree that
computes $f$, the maximum, over inputs $x$ with $f(x)=a$, of its expected number of
queries (with value $0$ assigned when $f(x)\ne a$) is nonnegative for every output
$a\in\alpha$.
\end{theorem}

\begin{theorem}[Input cost bounded by conditional cost]
\lean{FABL.RandomizedDecisionTree.inputCost_le_outputConditionalCost}
\label{FABL.RandomizedDecisionTree.inputCost_le_outputConditionalCost}
\leanok
\uses{FABL.RandomizedDecisionTree, FABL.RandomizedDecisionTree.inputCost, FABL.RandomizedDecisionTree.outputConditionalCost}
\difficulty{3}
Let $\Omega$ be a finite nonempty set, let $\alpha$ be a type with decidable equality,
let $n$ be a nonnegative integer, and let $f:(\mathrm{Fin} n\to\Omega)\to\alpha$. For a
zero-error randomized decision tree $T$ computing $f$, an output value $a\in\alpha$, and
an input $x$ satisfying $f(x)=a$, the expected number of queries on $x$ is at most the
maximum of that expected query count over all inputs on which $f$ equals $a$, with the
count taken to be zero on inputs producing another value.
\end{theorem}

\begin{definition}[False-output cost for recursive NAND]
\lean{FABL.recursiveNandFalseOutputCost}
\label{FABL.recursiveNandFalseOutputCost}
\leanok
\uses{FABL.RandomizedDecisionTree.outputConditionalCost, FABL.recursiveNandRandomizedTree}
For each depth $d$, define the false-output cost of the recursive randomized NAND tree
to be its output-conditional cost at output $1$.
\end{definition}

\begin{definition}[True-output cost for recursive NAND]
\lean{FABL.recursiveNandTrueOutputCost}
\label{FABL.recursiveNandTrueOutputCost}
\leanok
\uses{FABL.RandomizedDecisionTree.outputConditionalCost, FABL.recursiveNandRandomizedTree}
For each depth $d$, define the true-output cost of the recursive randomized NAND tree to
be its output-conditional cost at output $-1$.
\end{definition}

\begin{theorem}[Base-level randomized input cost]
\lean{FABL.recursiveNandRandomizedTree_inputCost_zero}
\label{FABL.recursiveNandRandomizedTree_inputCost_zero}
\leanok
\uses{FABL.FiniteDecisionTree.queryCount, FABL.FiniteDecisionTree.queryTrace, FABL.RandomizedDecisionTree.inputCost, FABL.SignCube, FABL.pmfExpectation, FABL.pmfExpectation_pure, FABL.recursiveNandIdentityTree, FABL.recursiveNandRandomizedTree, FABL.recursiveNandSeedLaw, FABL.recursiveNandTree}
\difficulty{3}
For every input $x\in\{\pm1\}^{2^0}$, the expected query cost of the level-zero
randomized tree is $1$.
\end{theorem}

\begin{theorem}[Initial false-output cost]
\lean{FABL.recursiveNandFalseOutputCost_zero}
\label{FABL.recursiveNandFalseOutputCost_zero}
\leanok
\uses{FABL.RandomizedDecisionTree.inputCost, FABL.RandomizedDecisionTree.outputConditionalCost, FABL.SignCube, FABL.recursiveNAND, FABL.recursiveNandFalseOutputCost, FABL.recursiveNandRandomizedTree, FABL.recursiveNandRandomizedTree_inputCost_zero}
\difficulty{4}
The recursive randomized decision tree at level $0$ has maximum expected query cost $1$
among inputs on which its output is $+1$.
\end{theorem}

\begin{theorem}[Initial true-output cost]
\lean{FABL.recursiveNandTrueOutputCost_zero}
\label{FABL.recursiveNandTrueOutputCost_zero}
\leanok
\uses{FABL.RandomizedDecisionTree.inputCost, FABL.RandomizedDecisionTree.outputConditionalCost, FABL.SignCube, FABL.recursiveNAND, FABL.recursiveNandRandomizedTree, FABL.recursiveNandRandomizedTree_inputCost_zero, FABL.recursiveNandTrueOutputCost}
\difficulty{4}
For the recursively defined randomized NAND decision tree, define the cost for an output
value as the maximum, over inputs producing that value, of the expected number of
queries. At depth zero, the cost for output $-1$ is $1$.
\end{theorem}

\begin{theorem}[Successor bound for false-output cost]
\lean{FABL.recursiveNandFalseOutputCost_succ_le}
\label{FABL.recursiveNandFalseOutputCost_succ_le}
\leanok
\uses{FABL.RandomizedDecisionTree.inputCost_le_outputConditionalCost, FABL.RandomizedDecisionTree.outputConditionalCost, FABL.RandomizedDecisionTree.outputConditionalCost_nonneg, FABL.nandFunction, FABL.recursiveInputBlock, FABL.recursiveNAND, FABL.recursiveNandFalseOutputCost, FABL.recursiveNandRandomizedTree, FABL.recursiveNandRandomizedTree_inputCost_succ, FABL.recursiveNandTrueOutputCost}
\difficulty{5}
For every natural number $d$, the maximum expected query cost on inputs with output $+1$
for the recursive randomized decision tree at level $d+1$ is at most twice the
corresponding maximum for output $-1$ at level $d$.
\end{theorem}

\begin{theorem}[Recursive true-output cost bound]
\lean{FABL.recursiveNandTrueOutputCost_succ_le}
\label{FABL.recursiveNandTrueOutputCost_succ_le}
\leanok
\uses{FABL.RandomizedDecisionTree.inputCost_le_outputConditionalCost, FABL.RandomizedDecisionTree.outputConditionalCost, FABL.RandomizedDecisionTree.outputConditionalCost_nonneg, FABL.nandFunction, FABL.recursiveInputBlock, FABL.recursiveNAND, FABL.recursiveNandFalseOutputCost, FABL.recursiveNandRandomizedTree, FABL.recursiveNandRandomizedTree_inputCost_succ, FABL.recursiveNandTrueOutputCost}
\difficulty{5}
For every natural number $d$, the recursively defined randomized NAND decision tree has
cost for output $-1$ at depth $d+1$ at most its cost for output $1$ at depth $d$ plus
half its cost for output $-1$ at depth $d$.
\end{theorem}

\begin{theorem}[Output-cost recurrence bounds]
\lean{FABL.recursiveNandOutputCosts_le_recurrence}
\label{FABL.recursiveNandOutputCosts_le_recurrence}
\leanok
\uses{FABL.recursiveNandFalseCost, FABL.recursiveNandFalseCost_succ, FABL.recursiveNandFalseOutputCost, FABL.recursiveNandFalseOutputCost_succ_le, FABL.recursiveNandTrueCost, FABL.recursiveNandTrueCost_succ, FABL.recursiveNandTrueOutputCost, FABL.recursiveNandTrueOutputCost_succ_le}
\difficulty{4}
For every natural number $d$, the randomized recursive NAND tree's maximum expected
query cost on inputs with output $+1$ is at most the recursively defined false cost at
level $d$, and its corresponding maximum on inputs with output $-1$ is at most the
recursively defined true cost at level $d$.
\end{theorem}

\begin{theorem}[Worst-case cost from output-conditioned costs]
\lean{FABL.recursiveNandRandomizedTree_worstCaseCost_le_maxOutputCost}
\label{FABL.recursiveNandRandomizedTree_worstCaseCost_le_maxOutputCost}
\leanok
\uses{FABL.RandomizedDecisionTree.inputCost_le_outputConditionalCost, FABL.RandomizedDecisionTree.worstCaseCost, FABL.recursiveNAND, FABL.recursiveNandFalseOutputCost, FABL.recursiveNandRandomizedTree, FABL.recursiveNandTrueOutputCost}
\difficulty{4}
For every $d\in\mathbb{N}$, the worst-case expected query cost of the level-$d$
randomized tree is at most the larger of its two output-conditioned costs: the supremum
of the input cost over inputs where the recursive NAND function equals $1$, with cost
taken to be zero on all other inputs, and the corresponding supremum for inputs where
the function equals $-1$.
\end{theorem}

\begin{theorem}[Worst-case cost recurrence bound]
\lean{FABL.recursiveNandRandomizedTree_worstCaseCost_le_recurrence}
\label{FABL.recursiveNandRandomizedTree_worstCaseCost_le_recurrence}
\leanok
\uses{FABL.RandomizedDecisionTree.worstCaseCost, FABL.recursiveNandFalseCost, FABL.recursiveNandOutputCosts_le_recurrence, FABL.recursiveNandRandomizedTree, FABL.recursiveNandRandomizedTree_worstCaseCost_le_maxOutputCost, FABL.recursiveNandTrueCost}
\difficulty{1}
For every $d\in\mathbb{N}$, the worst-case expected query cost of the level-$d$
randomized tree is at most $\max(F_d,T_d)$, where the real sequences $F_d$ and $T_d$ are
defined by $F_0=T_0=1$, $F_{d+1}=2T_d$, and $T_{d+1}=F_d+T_d/2$.
\end{theorem}

\begin{theorem}[Worst-case cost growth bound]
\lean{FABL.recursiveNandRandomizedTree_worstCaseCost_le_growth}
\label{FABL.recursiveNandRandomizedTree_worstCaseCost_le_growth}
\leanok
\uses{FABL.RandomizedDecisionTree.worstCaseCost, FABL.nandGrowthRate, FABL.recursiveNandRandomizedTree, FABL.recursiveNandRandomizedTree_worstCaseCost_le_recurrence, FABL.recursiveNand_maxCost_le}
\difficulty{1}
For every $d\in\mathbb{N}$, the worst-case expected query cost of the level-$d$
randomized tree is at most $\frac{2}{\rho}\rho^d$, where $\rho=(1+\sqrt{33})/4$.
\end{theorem}

\begin{theorem}[Randomized decision tree complexity recursive nand le]
\lean{FABL.randomizedDecisionTreeComplexity_recursiveNAND_le}
\statementsource{Analysis-of-Boolean-Functions-by-Ryan-ODonnell}{pdf-ffbe9c7b3306-p242-b005}
\label{FABL.randomizedDecisionTreeComplexity_recursiveNAND_le}
\leanok
\uses{FABL.RandomizedDecisionTree, FABL.RandomizedDecisionTree.complexity, FABL.RandomizedDecisionTree.worstCaseCost, FABL.RandomizedDecisionTree.worstCaseCost_nonneg, FABL.Sign, FABL.nandGrowthRate, FABL.recursiveNAND, FABL.recursiveNandRandomizedTree, FABL.recursiveNandRandomizedTree_worstCaseCost_le_growth}
\difficulty{4}
% dataset-pipeline: stub -- prose pending build_informal_statements.py
Statement pending regeneration.
\end{theorem}

\begin{theorem}[False-output cost at level one]
\lean{FABL.recursiveNandFalseOutputCost_one}
\label{FABL.recursiveNandFalseOutputCost_one}
\leanok
\uses{FABL.RandomizedDecisionTree.inputCost, FABL.RandomizedDecisionTree.inputCost_le_outputConditionalCost, FABL.SignCube, FABL.nandFunction, FABL.recursiveInputBlock, FABL.recursiveNAND, FABL.recursiveNandFalseCost, FABL.recursiveNandFalseOutputCost, FABL.recursiveNandOutputCosts_le_recurrence, FABL.recursiveNandRandomizedTree, FABL.recursiveNandRandomizedTree_inputCost_succ, FABL.recursiveNandRandomizedTree_inputCost_zero, FABL.recursiveNand_costs_one}
\difficulty{4}
For the recursively defined randomized decision tree on two sign inputs, the maximum
expected number of queries over inputs whose output is $+1$ is $2$.
\end{theorem}

\begin{theorem}[True-output cost at depth one]
\lean{FABL.recursiveNandTrueOutputCost_one}
\label{FABL.recursiveNandTrueOutputCost_one}
\leanok
\uses{FABL.RandomizedDecisionTree.inputCost, FABL.RandomizedDecisionTree.inputCost_le_outputConditionalCost, FABL.SignCube, FABL.nandFunction, FABL.recursiveInputBlock, FABL.recursiveNAND, FABL.recursiveNandOutputCosts_le_recurrence, FABL.recursiveNandRandomizedTree, FABL.recursiveNandRandomizedTree_inputCost_succ, FABL.recursiveNandRandomizedTree_inputCost_zero, FABL.recursiveNandTrueCost, FABL.recursiveNandTrueOutputCost, FABL.recursiveNand_costs_one}
\difficulty{5}
For the recursively defined randomized NAND decision tree, define the cost for an output
value as the maximum, over inputs producing that value, of the expected number of
queries. At depth one, the cost for output $-1$ is $3/2$.
\end{theorem}

\begin{theorem}[Depth-one false-output cost differs]
\lean{FABL.exercise8_40d_printed_false_label_is_wrong}
\label{FABL.exercise8_40d_printed_false_label_is_wrong}
\leanok
\uses{FABL.recursiveNandFalseOutputCost, FABL.recursiveNandFalseOutputCost_one}
\difficulty{1}
For the randomized decision tree defined by the recursive NAND construction at depth
one, its output-conditional cost for output $1$ is not equal to $3/2$.
\end{theorem}

\begin{theorem}[Level-one worst-case cost]
\lean{FABL.recursiveNandRandomizedTree_worstCaseCost_one}
\label{FABL.recursiveNandRandomizedTree_worstCaseCost_one}
\leanok
\uses{FABL.RandomizedDecisionTree.inputCost, FABL.RandomizedDecisionTree.inputCost_le_worstCaseCost, FABL.RandomizedDecisionTree.worstCaseCost, FABL.SignCube, FABL.recursiveInputBlock, FABL.recursiveNAND, FABL.recursiveNandFalseOutputCost, FABL.recursiveNandFalseOutputCost_one, FABL.recursiveNandRandomizedTree, FABL.recursiveNandRandomizedTree_inputCost_succ, FABL.recursiveNandRandomizedTree_inputCost_zero, FABL.recursiveNandRandomizedTree_worstCaseCost_le_maxOutputCost, FABL.recursiveNandTrueOutputCost, FABL.recursiveNandTrueOutputCost_one}
\difficulty{4}
The worst-case expected query cost of the level-one randomized tree is $2$.
\end{theorem}

\begin{theorem}[NAND growth rate is below two]
\lean{FABL.nandGrowthRate_lt_two}
\label{FABL.nandGrowthRate_lt_two}
\leanok
\uses{FABL.nandGrowthRate}
\difficulty{2}
The real number $(1+\sqrt{33})/4$ is strictly less than $2$.
\end{theorem}

\begin{theorem}[Depth-one cost exceeds the growth rate]
\lean{FABL.exercise8_40g_coefficient_one_fails_at_depth_one}
\label{FABL.exercise8_40g_coefficient_one_fails_at_depth_one}
\leanok
\uses{FABL.RandomizedDecisionTree.worstCaseCost, FABL.nandGrowthRate, FABL.nandGrowthRate_lt_two, FABL.recursiveNandRandomizedTree, FABL.recursiveNandRandomizedTree_worstCaseCost_one}
\difficulty{1}
For the recursive randomized decision tree at depth one, its worst-case expected number
of queries is not at most $\lambda$, where $\lambda=(1+\sqrt{33})/4$.
\end{theorem}

\begin{definition}[Formula address space]
\lean{FABL.FormulaAddress}
\label{FABL.FormulaAddress}
\leanok
For an arity $r$, define the address space at depth $0$ to be a singleton. At depth
$d+1$, an address consists of a choice from $\{0,\ldots,r-1\}$ followed by an address of
depth $d$.
\end{definition}

\begin{definition}[Finite formula address equivalence]
\lean{FABL.finFormulaAddressEquiv}
\label{FABL.finFormulaAddressEquiv}
\leanok
\uses{FABL.FormulaAddress}
For each arity $r$, define recursively an equivalence between the integers
$0,\ldots,r^d-1$ and the addresses of depth $d$ in an $r$-ary formula.
\end{definition}

\begin{definition}[Recursive binary NAND evaluation]
\lean{FABL.addressRecursiveNAND}
\label{FABL.addressRecursiveNAND}
\leanok
\uses{FABL.FormulaAddress, FABL.Sign, FABL.nandPair}
For each depth $d$, define the value of a binary NAND formula whose leaves are indexed
by depth-$d$ addresses: at depth zero, return the value at the unique address; at
positive depth, take the NAND of the recursively evaluated subformulas indexed by first
digit $0$ and first digit $1$.
\end{definition}

\begin{theorem}[Recursive NAND on addressed inputs]
\lean{FABL.recursiveNAND_eq_addressRecursiveNAND}
\label{FABL.recursiveNAND_eq_addressRecursiveNAND}
\leanok
\uses{FABL.SignCube, FABL.addressRecursiveNAND, FABL.finFormulaAddressEquiv, FABL.nandPair, FABL.recursiveInputBlock, FABL.recursiveNAND, FABL.recursiveNAND_succ}
\difficulty{4}
For every natural number $d$ and every sign assignment $x$ to the $2^d$ coordinates, the
depth-$d$ recursive NAND equals the depth-$d$ address-recursive NAND on the assignment
that gives each formula address the value of $x$ at the corresponding coordinate under
the canonical bijection between the $2^d$ coordinates and depth-$d$ formula addresses of
arity two.
\end{theorem}

\begin{definition}[Reindexing sign-valued functions]
\lean{FABL.piReindexEquiv}
\label{FABL.piReindexEquiv}
\leanok
\uses{FABL.Sign}
Given types $\iota$ and $\kappa$ and a bijection $e:\iota\to\kappa$, the induced
bijection between sign-valued functions on these types sends $x$ to the function
$k\mapsto x(e^{-1}(k))$.
\end{definition}

\begin{definition}[Top correlation]
\lean{FABL.topCorrelation}
\label{FABL.topCorrelation}
\leanok
\uses{FABL.Sign, FABL.signValue}
For a finite index set $\iota$ and a real-valued function $f$ on sign-valued assignments
$x:\iota\to\{\pm1\}$, define its top correlation to be the uniform average of
$f(x)\prod_{i\in\iota}x(i)$.
\end{definition}

\begin{theorem}[Reindexing top correlation]
\lean{FABL.topCorrelation_reindex}
\label{FABL.topCorrelation_reindex}
\leanok
\uses{FABL.Sign, FABL.piReindexEquiv, FABL.topCorrelation}
\difficulty{3}
Let $I$ and $K$ be finite index sets, let $e:I\to K$ be a bijection, and let $f$ be a
real-valued function on sign assignments indexed by $K$. The average of $f$ times the
product of all coordinate signs is unchanged when the assignments are reindexed along
$e$.
\end{theorem}

\begin{theorem}[Mean of full parity]
\lean{FABL.expect_fullParity_eq_zero}
\label{FABL.expect_fullParity_eq_zero}
\leanok
\uses{FABL.Sign, FABL.SignCube, FABL.expect_monomial, FABL.monomial, FABL.piReindexEquiv, FABL.signValue}
\difficulty{5}
Let $\iota$ be a finite, nonempty set, and choose independent uniform signs
$x_i\in\{-1,1\}$ for each $i\in\iota$. The expected value of their product is zero:
$\mathbb E\prod_{i\in\iota}x_i=0$.
\end{theorem}

\begin{definition}[Top coefficient of recursive NAND]
\lean{FABL.addressRecursiveNANDTopCoeff}
\label{FABL.addressRecursiveNANDTopCoeff}
\leanok
\uses{FABL.addressRecursiveNAND, FABL.signValue, FABL.topCorrelation}
For each natural-number depth $d$, define the top coefficient of the recursively
evaluated binary NAND formula as the top correlation of the function that maps each
assignment of signs to the real sign value of the formula's output.
\end{definition}

\begin{theorem}[Top coefficient of recursive NAND]
\lean{FABL.fourierCoeff_recursiveNAND_univ_eq_addressTopCoeff}
\label{FABL.fourierCoeff_recursiveNAND_univ_eq_addressTopCoeff}
\leanok
\uses{FABL.BooleanFunction.toReal, FABL.addressRecursiveNAND, FABL.addressRecursiveNANDTopCoeff, FABL.finFormulaAddressEquiv, FABL.fourierCoeff, FABL.monomial, FABL.piReindexEquiv, FABL.recursiveNAND, FABL.recursiveNAND_eq_addressRecursiveNAND, FABL.signValue, FABL.topCorrelation, FABL.topCorrelation_reindex}
\difficulty{3}
For every nonnegative integer $d$, the Fourier coefficient on all coordinates of the
depth-$d$ recursive NAND function equals the top correlation of the corresponding
recursive NAND formula on its leaf addresses.
\end{theorem}

\begin{definition}[Splitting two blocks of coordinates]
\lean{FABL.splitTwoBlocksEquiv}
\label{FABL.splitTwoBlocksEquiv}
\leanok
\uses{FABL.Sign}
For any type $A$, there is a bijection from sign-valued functions on
$\mathrm{Fin}(2)\times A$ to pairs of sign-valued functions on $A$. It restricts a
function to the block with first coordinate $0$ and to the block with first coordinate
$1$; the inverse combines the two blocks according to that first coordinate.
\end{definition}

\begin{theorem}[Value of a NAND sign pair]
\lean{FABL.signValue_nandPair}
\label{FABL.signValue_nandPair}
\leanok
\uses{FABL.Sign, FABL.nandPair, FABL.signValue}
\difficulty{4}
For any signs $a$ and $b$, the NAND sign pair has real value $\frac{-1-a-b+ab}{2}$.
\end{theorem}

\begin{theorem}[Factorization of a product expectation]
\lean{FABL.expect_pair_mul_separate}
\label{FABL.expect_pair_mul_separate}
\leanok
\difficulty{3}
Let $A$ and $B$ be finite sets, and let $f\colon A\to\mathbb R$ and $g\colon B\to\mathbb
R$. Under the uniform distribution on $A\times B$, the expected value of $f(a)g(b)$ is
the product of the uniform expectations of $f$ on $A$ and $g$ on $B$.
\end{theorem}

\begin{theorem}[Expectation of a scaled function]
\lean{FABL.expect_div_const}
\label{FABL.expect_div_const}
\leanok
\difficulty{3}
Let $A$ be a finite set, let $f:A\to\mathbb R$, and let $c\in\mathbb R$. The uniform
average of $a\mapsto f(a)/c$ equals the uniform average of $f$ divided by $c$.
\end{theorem}

\begin{theorem}[Averaged NAND parity identity]
\lean{FABL.expect_nandParity_step}
\label{FABL.expect_nandParity_step}
\leanok
\uses{FABL.expect_div_const, FABL.expect_pair_mul_separate}
\difficulty{5}
Let $A$ be a finite nonempty set, and let $F,P:A\to\mathbb R$ satisfy
$|A|^{-1}\sum_{a\in A}P(a)=0$. Then
\[
\frac{1}{|A|^2}\sum_{a,b\in A}\frac{-1-F(a)-F(b)+F(a)F(b)}{2}P(a)P(b)
=\frac12\left(\frac1{|A|}\sum_{a\in A}F(a)P(a)\right)^2.
\]
\end{theorem}

\begin{theorem}[Initial NAND top coefficient]
\lean{FABL.addressRecursiveNANDTopCoeff_zero}
\label{FABL.addressRecursiveNANDTopCoeff_zero}
\leanok
\uses{FABL.Sign, FABL.addressRecursiveNAND, FABL.addressRecursiveNANDTopCoeff, FABL.signValue, FABL.topCorrelation}
\difficulty{4}
The top coefficient of the depth-zero recursively composed two-input NAND function is
$1$.
\end{theorem}

\begin{theorem}[Recurrence for the NAND top coefficient]
\lean{FABL.addressRecursiveNANDTopCoeff_succ}
\label{FABL.addressRecursiveNANDTopCoeff_succ}
\leanok
\uses{FABL.FormulaAddress, FABL.Sign, FABL.addressRecursiveNAND, FABL.addressRecursiveNANDTopCoeff, FABL.expect_fullParity_eq_zero, FABL.expect_nandParity_step, FABL.nandPair, FABL.signValue, FABL.signValue_nandPair, FABL.splitTwoBlocksEquiv, FABL.topCorrelation}
\difficulty{4}
For every depth $d\in\mathbb{N}$, the top coefficient of the depth-$(d+1)$ recursively
composed two-input NAND function equals one half the square of the depth-$d$ top
coefficient.
\end{theorem}

\begin{theorem}[Nonvanishing of the NAND top coefficient]
\lean{FABL.addressRecursiveNANDTopCoeff_ne_zero}
\label{FABL.addressRecursiveNANDTopCoeff_ne_zero}
\leanok
\uses{FABL.addressRecursiveNANDTopCoeff, FABL.addressRecursiveNANDTopCoeff_succ}
\difficulty{4}
For every depth $d\in\mathbb{N}$, the top correlation of the depth-$d$ recursively
composed two-input NAND function is nonzero.
\end{theorem}

\begin{theorem}[Nonvanishing of the recursive NAND coefficient]
\lean{FABL.fourierCoeff_recursiveNAND_univ_ne_zero}
\label{FABL.fourierCoeff_recursiveNAND_univ_ne_zero}
\leanok
\uses{FABL.BooleanFunction.toReal, FABL.addressRecursiveNANDTopCoeff_ne_zero, FABL.fourierCoeff, FABL.fourierCoeff_recursiveNAND_univ_eq_addressTopCoeff, FABL.recursiveNAND}
\difficulty{3}
For every nonnegative integer $d$, the Fourier coefficient on all coordinates of the
depth-$d$ recursive NAND function is nonzero.
\end{theorem}

\begin{theorem}[Degree of recursive NAND]
\lean{FABL.fourierDegree_recursiveNAND}
\label{FABL.fourierDegree_recursiveNAND}
\leanok
\uses{FABL.BooleanFunction.toReal, FABL.fourierCoeff_recursiveNAND_univ_ne_zero, FABL.fourierDegree, FABL.fourierDegree_le_dimension, FABL.fourierDegree_le_iff, FABL.recursiveNAND}
\difficulty{4}
For every natural number $d$, the real-valued function obtained by converting the
recursively defined NAND function of depth $d$ to real values has Fourier degree exactly
$2^d$.
\end{theorem}

\begin{definition}[Conversion to a binary real decision tree]
\lean{FABL.FiniteDecisionTree.toF₂Real}
\leanok
\uses{FABL.FiniteDecisionTree, FABL.F₂DecisionTree, FABL.Sign, FABL.signValue}
Given a finite decision tree with sign-valued labels and sign-valued queries, indexed by
a set of available coordinates among $\mathrm{Fin}(n)$, define a binary decision tree
with real labels on the same available coordinates. Each leaf label is converted to its
real sign value; each query is retained, with its zero child taken from the original
child indexed by $1$ and its one child from the child indexed by $-1$.
\end{definition}

\begin{theorem}[Evaluation under binary sign encoding]
\lean{FABL.FiniteDecisionTree.eval_toF₂Real}
\leanok
\uses{FABL.FiniteDecisionTree, FABL.FiniteDecisionTree.eval, FABL.FiniteDecisionTree.toF₂Real, FABL.F₂Cube, FABL.F₂DecisionTree.eval, FABL.Sign, FABL.binaryCubeSignEquiv, FABL.binaryCubeSignEquiv_apply, FABL.binarySignEquiv, FABL.signEncode, FABL.signEncode_zero, FABL.signValue}
\difficulty{5}
Let $T$ be a finite decision tree on $n$ coordinates whose inputs and outputs are signs
in $\{ -1,1\}$, with available coordinates $A$, and let $x$ be a binary vector in
$(\mathbb{Z}/2\mathbb{Z})^n$. The associated $\mathbb{Z}/2\mathbb{Z}$-valued decision
tree evaluated at $x$ equals the real value of the sign obtained by evaluating $T$ at
the coordinatewise sign encoding of $x$.
\end{theorem}

\begin{theorem}[Depth under conversion to a binary tree]
\lean{FABL.FiniteDecisionTree.depth_toF₂Real_le}
\leanok
\uses{FABL.FiniteDecisionTree, FABL.FiniteDecisionTree.depth, FABL.FiniteDecisionTree.toF₂Real, FABL.F₂DecisionTree.depth, FABL.Sign}
\difficulty{4}
Let $A$ be a set of available coordinates among $n$ coordinates, and let $T$ be a
decision tree with sign-valued inputs and outputs over $A$. The associated binary
decision tree has depth at most the depth of $T$.
\end{theorem}

\begin{theorem}[Real-valued conversion preserves computation]
\lean{FABL.FiniteDecisionTree.toF₂Real_computes}
\leanok
\uses{FABL.BooleanFunction, FABL.BooleanFunction.toReal, FABL.FiniteDecisionTree, FABL.FiniteDecisionTree.Computes, FABL.FiniteDecisionTree.eval_toF₂Real, FABL.FiniteDecisionTree.toF₂Real, FABL.F₂DecisionTree.Computes, FABL.Sign, FABL.binaryCubeSignEquiv}
\difficulty{1}
Let $n$ be a natural number, let $A$ be any subset of the $n$ coordinates, let $f$ map
the sign cube of dimension $n$ to the two signs, and let $T$ be a finite decision tree
with sign-valued queries and outputs whose available coordinates are in $A$. If $T$
computes $f$, then the corresponding binary decision tree computes the real-valued
function obtained by converting each binary input coordinate to its sign and
interpreting the output sign of $f$ as a real number.
\end{theorem}

\begin{theorem}[Depth lower bound for recursive NAND]
\lean{FABL.FiniteDecisionTree.recursiveNAND_depth_le}
\label{FABL.FiniteDecisionTree.recursiveNAND_depth_le}
\leanok
\uses{FABL.FiniteAlphabetDecisionTree, FABL.FiniteDecisionTree.Computes, FABL.FiniteDecisionTree.depth, FABL.FiniteDecisionTree.depth_toF₂Real_le, FABL.FiniteDecisionTree.toF₂Real, FABL.FiniteDecisionTree.toF₂Real_computes, FABL.Sign, FABL.fourierDegree_recursiveNAND, FABL.fourierDegree_toReal_le_depth_of_decisionTree, FABL.recursiveNAND}
\difficulty{3}
For every natural number $d$, any finite decision tree on the $2^d$-dimensional sign
cube that computes the recursively defined NAND function has depth at least $2^d$.
\end{theorem}

\begin{theorem}[Decision-tree depth of recursive NAND]
\lean{FABL.deterministicDepth_recursiveNAND}
\label{FABL.deterministicDepth_recursiveNAND}
\leanok
\uses{FABL.FiniteDecisionTree.deterministicDepth, FABL.FiniteDecisionTree.exists_computingTree_depth_eq_deterministicDepth, FABL.FiniteDecisionTree.recursiveNAND_depth_le, FABL.deterministicDepth_recursiveNAND_le, FABL.recursiveNAND}
\difficulty{4}
For every nonnegative integer $d$, define a Boolean function on $2^d$ sign-valued inputs
recursively: at depth zero it returns its single input, and each subsequent level
applies the two-input NAND function to the outputs on the two recursively defined input
blocks. The minimum depth of a deterministic decision tree computing this function is
$2^d$.
\end{theorem}

\begin{definition}[Pairing four symbols as two binary symbols]
\lean{FABL.finFourPairEquiv}
\label{FABL.finFourPairEquiv}
\leanok
There is a bijection between $\mathrm{Fin}(4)$ and
$\mathrm{Fin}(2)\times\mathrm{Fin}(2)$, identifying each of four symbols with an ordered
pair of binary symbols.
\end{definition}

\begin{definition}[Even-depth formula address bijection]
\lean{FABL.evenFormulaAddressEquiv}
\label{FABL.evenFormulaAddressEquiv}
\leanok
\uses{FABL.FormulaAddress, FABL.finFourPairEquiv}
For every natural number $k$, there is a bijection between the addresses of a $4$-ary
formula of depth $k$ and the addresses of a binary formula of depth $2k$.
\end{definition}

\begin{definition}[Addressed iterated tribes]
\lean{FABL.addressIteratedTribesTwoTwo}
\label{FABL.addressIteratedTribesTwoTwo}
\leanok
\uses{FABL.FormulaAddress, FABL.Sign, FABL.tribes}
For each natural number $k$, define a Boolean function on assignments of signs to the
depth-$k$ addresses of a quaternary tree as follows: at depth zero, return the sign at
the unique address; at depth $k+1$, group the addresses by their first coordinate, apply
the depth-$k$ function to each resulting subtree, and apply the tribes function with
block width and number of blocks both equal to $2$ to the two outputs.
\end{definition}

\begin{theorem}[Two-level NAND identity for tribes]
\lean{FABL.nandPair_two_levels_eq_tribes}
\label{FABL.nandPair_two_levels_eq_tribes}
\leanok
\uses{FABL.SignCube, FABL.andFunction, FABL.finFourPairEquiv, FABL.inputBlock, FABL.nandPair, FABL.orFunction, FABL.tribes}
\difficulty{5}
For every $x\in\{-1,1\}^4$, apply the two-input operation that returns $1$ exactly when
both entries are $-1$, and $-1$ otherwise, first to each of the two coordinate pairs and
then to those two outputs. The result is $1$ exactly when each coordinate pair contains
at least one $1$, and is $-1$ otherwise.
\end{theorem}

\begin{theorem}[Even-depth NAND and iterated tribes]
\lean{FABL.addressRecursiveNAND_even_eq_iteratedTribes}
\label{FABL.addressRecursiveNAND_even_eq_iteratedTribes}
\leanok
\uses{FABL.FormulaAddress, FABL.Sign, FABL.addressIteratedTribesTwoTwo, FABL.addressRecursiveNAND, FABL.evenFormulaAddressEquiv, FABL.finFourPairEquiv, FABL.nandPair, FABL.nandPair_two_levels_eq_tribes}
\difficulty{5}
For every $k\in\mathbb{N}$ and every assignment $x$ of signs to the quaternary addresses
of depth $k$, the depth-$2k$ recursively composed two-input NAND function, given the
input at each binary address by reading $x$ at the corresponding quaternary address
under the canonical bijection that pairs successive binary digits, equals the depth-$k$
iterated two-by-two tribes function on $x$.
\end{theorem}

\begin{definition}[Iterated tribes on the Boolean cube]
\lean{FABL.iteratedTribesTwoTwo}
\label{FABL.iteratedTribesTwoTwo}
\leanok
\uses{FABL.BooleanFunction, FABL.addressIteratedTribesTwoTwo, FABL.finFormulaAddressEquiv}
For each natural number $k$, define a Boolean function on the sign cube of dimension
$4^k$ by identifying its coordinates with the depth-$k$ addresses of a quaternary tree
and applying the addressed iterated tribes function.
\end{definition}

\begin{definition}[Recursive NAND coordinate equivalence]
\lean{FABL.recursiveNandEvenTribesCoordinateEquiv}
\label{FABL.recursiveNandEvenTribesCoordinateEquiv}
\leanok
\uses{FABL.evenFormulaAddressEquiv, FABL.finFormulaAddressEquiv}
For each natural number $k$, the composite of the depth-$2k$ binary-address coordinate
equivalence and the inverse address correspondences between even binary addresses and
quaternary addresses gives an equivalence from the coordinates indexed by
$	ext{Fin}(2^{2k})$ to those indexed by $	ext{Fin}(4^k)$.
\end{definition}

\begin{theorem}[Even recursive NAND as iterated tribes]
\lean{FABL.recursiveNAND_even_eq_iteratedTribesTwoTwo}
\label{FABL.recursiveNAND_even_eq_iteratedTribesTwoTwo}
\leanok
\uses{FABL.FormulaAddress, FABL.Sign, FABL.SignCube, FABL.addressRecursiveNAND_even_eq_iteratedTribes, FABL.evenFormulaAddressEquiv, FABL.finFormulaAddressEquiv, FABL.iteratedTribesTwoTwo, FABL.piReindexEquiv, FABL.recursiveNAND, FABL.recursiveNAND_eq_addressRecursiveNAND, FABL.recursiveNandEvenTribesCoordinateEquiv}
\difficulty{3}
For every natural number $k$ and every sign assignment $x$ to the $2^{2k}$ coordinates,
the recursive NAND of depth $2k$ equals the iterated tribes function of depth $k$ on the
assignment reindexed by the canonical bijection from the $2^{2k}$ coordinates to the
$4^k$ coordinates.
\end{theorem}
